\documentclass[11pt,a4paper]{article}

\usepackage[utf8]{inputenc}
\usepackage[T1]{fontenc}
\usepackage[spanish]{babel}
\usepackage[margin=2.5cm]{geometry}
\usepackage{mathptmx}
\usepackage{enumitem}
\usepackage[hidelinks]{hyperref}
\usepackage{microtype}

\setlength{\parindent}{0pt}
\setlength{\parskip}{6pt}
\setlist[itemize]{leftmargin=*, itemsep=2pt, topsep=4pt}
\pagestyle{empty}

\begin{document}

\begin{center}
    {\large\bfseries PONTIFICIA UNIVERSIDAD CATÓLICA DE VALPARAÍSO}\\[0.3em]
    Facultad de Ingeniería\\[0.2em]
    Escuela de Ingeniería Informática\\[1.2em]
    {\Large\bfseries TRABAJO FINAL}\\[0.4em]
    INF3541-2 \quad TALLER DE BASE DE DATOS
\end{center}

\vspace{0.4em}
\hrule
\vspace{1em}

Como finalización de Ramo, deberán entregar un trabajo realizado en el Sistema
Administrador de Bases de Datos PostgresSQL el cual consiste en implementar una
base de datos a partir de un gran conjunto de datos, según elección en listado
Adjunto.

\begin{itemize}
    \item Seleccione un conjunto de datos con una cantidad considerables de
    registros, que al normalizarlo obtenga alrededor de 10 tablas y que sea
    reciente (año 2025 o 2026), desde alguno de estos sitios:
    \begin{itemize}
        \item \url{https://www.kaggle.com/datasets}
        \item \url{https://datahub.io/collections}
        \item Búsqueda por temas en
        \url{https://datasetsearch.research.google.com}
    \end{itemize}

    \item Modele y cargue el set de datos elegido. El modelo creado debe manejar
    control total de integridad y reglas del negocio totalmente controladas
    mediante el uso de \texttt{constraint} y/o \texttt{triggers}.

    \item Sobre el modelo cargado, genere:
    \begin{itemize}
        \item Uso de \texttt{triggers} para manejo de Logs de transacciones sobre
        tabla (s) principal.

        \item Generación de procedimientos almacenados para los mantenedores de
        las tablas. Al menos uno de ellos debe implementar alguna regla del
        negocio (validación, etc). Debe manejar correctamente las transacciones.

        \item Generación de consultas (informes) los cuales deben incluir
        consultas complejas (\texttt{subqueries}).
        \begin{itemize}
            \item Realice una consulta restringiendo los datos, que haga uso del
            operadores OR, operador IN, UNION o usando enlazando las tablas por
            el \texttt{where} (sin uso de la sentencia \texttt{join}).

            \item Realice una consulta equivalente a la anterior (que retorne los
            mismos resultados), pero que la sentencia sea distinta. Comprare
            ambos planes de ejecución, y trate de mejorarlos, ya sea modificando
            la consulta y/o creando índices.
        \end{itemize}

        \item Manejo de usuario con distinto nivel de acceso a los datos (incluso
        con filtros mediante vistas).
    \end{itemize}
\end{itemize}

\medskip

\noindent\textbf{Forma de entrega:} Se debe entregar un archivos ZIP que contenga
script de creación de la base de datos, la carga de los datos a tabla temporal y
el traspaso de los datos a tablas del modelo. Además entregar script de las
funciones / \texttt{triggers} / consultas, y en forma adicional un informe en PDF
muestre las tareas realizadas, los planes de ejecución previo y posterior a las
modificaciones, y las conclusiones de la optimización de consulta.

\medskip

\noindent\textbf{Fecha Máxima selección del tema:} martes 10 de noviembre.

\medskip

\noindent\textbf{Fecha Entrega final:} miércoles 25 de noviembre.

\end{document}
